% =========================================================
% PARTE 4: ALGORITMO UCB (MULTI-ARMED BANDIT)
% =========================================================
\section{Parte 4: Algoritmo UCB1 (Multi-Armed Bandit)}

Para superar la rigidez de la política tradicional de combi llena ($18$ pasajeros), se formula el problema de decisión de despacho como un entorno de \textbf{Aprendizaje por Refuerzo} utilizando el algoritmo \textbf{Upper Confidence Bound (UCB1)}.

\subsection{Definición de Brazos (Políticas de Despacho)}

El sistema selecciona en cada decisión entre 5 brazos o políticas candidatas ($a_0, \dots, a_4$):

\begin{table}[H]
	\centering
	\caption{Catálogo de Políticas de Salida (Brazos UCB)}
	\label{tab:brazos_catalog}
	\begin{tabularx}{\textwidth}{>{\bfseries\color{purpleUPP}}c l X}
		\toprule
		Brazo & Identificador & Regla de Despacho \\
		\midrule
		$a_0$ & \texttt{FullOnly} & Salir únicamente al alcanzar $18$ pasajeros (Política Tradicional). \\
		$a_1$ & \texttt{Flex\_14p\_15m} & Salir con $\ge 14$ personas O tras 15 minutos de espera. \\
		$a_2$ & \texttt{Flex\_12p\_25m} & Salir con $\ge 12$ personas O tras 25 minutos de espera. \\
		$a_3$ & \texttt{Flex\_10p\_35m} & Salir con $\ge 10$ personas O tras 35 minutos de espera. \\
		$a_4$ & \texttt{Flex\_8p\_45m} & Salir con $\ge 8$ personas O tras 45 minutos de espera. \\
		\bottomrule
	\end{tabularx}
\end{table}

\subsection{Función de Recompensa Normalizada}

En cada turno $t$, la penalización total del sistema se calcula ponderando tres métricas clave:
\begin{equation}
	P_t = w_1 \cdot \text{espera\_total} + w_2 \cdot \text{alumnos\_dejados} + w_3 \cdot \text{asientos\_vacíos}
\end{equation}

La recompensa inmediata $R_t \in [0, 1]$ se obtiene mediante min-max scaling invertido:
\begin{equation}
	R_t = 1 - \frac{P_t - P_{\text{min}}}{P_{\text{max}} - P_{\text{min}}}
\end{equation}

\subsection{Regla de Selección UCB1}

En el paso $t$, la IA selecciona la acción $a^*$ que maximiza el valor esperado más la cota de incertidumbre:
\begin{equation}
	a^* = \arg\max_{a \in \{a_0, \dots, a_4\}} \left[ \bar{Q}(a) + c \cdot \sqrt{\frac{\ln t}{N_a(t)}} \right]
\end{equation}

donde $\bar{Q}(a)$ es la recompensa promedio acumulada por el brazo $a$, $N_a(t)$ es la cantidad de veces que ha sido elegido, y $c = \sqrt{2}$ es el coeficiente de exploración.

\newpage
